\chapter{Použité technologie}
\label{3-technologie}

Třetí kapitola se zabývá užitými technologiemi, stručně představuje jejich historii a~zaměření.
Podrobnější informace o~jednotlivých projektech lze nalézt na oficiálních internetových stránkách. 


\section{Python}
\label{python}

  \begin{figure}[H]
   \centering
	\includegraphics[scale=0.5]{./pictures/python-logo-master-v3-TM.png}
	\caption[Python logo]{Python logo 
      (zdroj: \url{https://www.python.org/static/community_logos/python-logo-master-v3-TM.png})}
      \label{fig:python}
  \end{figure}

Python je open-source multiplatformní programovací interpretovaný jazyk pojme\-novaný podle
kultovního britského komediálního seriálu Monty Pythonův létající cirkus. Momentálně se
vyvíjejí verze třetí řady (3.X). 

První verzi, vydanou roku 1991, navrhl nizozemský programátor a~inženýr Guido van~Rossum za užití kódu
mnoha dalších programátorů. Silnou inspiraci van~Rossum našel v~jazyku ABC. Odtud také plyne
přívětivost jazyka, pro niž je často vyučován a~doporučován jako ideální pro rychlý
(a~přesto kvalitní) vývoj aplikací. 

Zásadní byl pro vývoj rok~2001, kdy vznikla nezisková organizace Python Software
Foundation (\zk{PSF}). \zk{PSF} zaštiťuje další vývoj Pythonu, vlastní jeho
\textit{intelektuální jmění} a~pořádá a~podporuje konference ohledně Pythonu. Sám van~Rossum bývá
pro svůj dohled a~svou neomezenou možnost zasahovat do vývoje nazýván \textit{benevolentní doživotní
diktátor}, španělská inkvizice jazyka Python. 

Python se na první pohled vyznačuje odsazováním kódu. Toho se ve většině jazyků používá jen jako
pomůcky pro přehlednost, v~Pythonu je však odsazování povinné. 

Mezi další specifikace jazyka patří například: 
\begin{itemize}

	\item Chování funkce: Ta se do svého zavolání uchovává jako objekt. 
	
	\item Proměnné: Není třeba deklarovat jejich typ. To může ušetřit značné množství práce
	v~psaní kódu, ale je třeba si na tuto vlastnost dát pozor (celé číslo dělené celým číslem je
	vždy celé číslo, i kdyby správným výsledkem mělo být číslo desetinné). 
	
	\item Proměnné uvnitř objektu: Není třeba je udávat při vytváření objektu. Lze je založit později. 
	
	\item Dokončení zápisu: V~interaktivním režimu provádím dokončení zápisu prázdným řádkem. 
	
	\item Kompilace: Python si automaticky kompiluje moduly do souborů \textit{.pyc}. Zkompilovaný
	modul je platformně nezávislý. 

\end{itemize}


\section{QGIS}
\label{qgis}

  \begin{figure}[H]
    \centering
      \includegraphics[width=120pt]{./pictures/qgis.png}
      \caption[QGIS logo]{QGIS logo 
      (zdroj: \href{https://commons.wikimedia.org/wiki/File:QGIS\_logo.svg}{Wikimedia Commons})}
      \label{fig:qgis}
  \end{figure}

QGIS (zkratka dříve užívaného názvu Quantum~GIS) představuje na poli geogra\-fic\-kých informačních systémů (\zk{GIS}) open-source alternativu ke komerčním geografic\-kým informačním systémům typu~ArcGIS. 

Na začátku 21.~století už byly geografické informační systémy běžnou praxí. Zrodila se tedy potřeba volby svobodného multiplatformního systému s~širokou podporou formátů geodat. Roku 2002 začal Quantum~GIS vyvíjet Gary~Sherman, k~němuž se později připojovalo čím dál více dobrovolníků. Časem projekt zaštítilo též Open Source Geospatial Foundation (\zk{OSGeo})~–~organizace pro podporu a~vývoj otevřených geoinformačních technologií a dat vzniklá roku~2006. Verze~1.0 byla uveřejněna v~lednu~2009. 

První verze systému QGIS byly pojmenovávány podle psů, později se přešlo na jména měsíců Jupiteru a~Saturnu. Momentálně nesou nejnovější verze názvy měst. 
QGIS se ve svém zaměření příliš neliší od běžných geografických informačních systémů. Uživatel v~něm má možnost prohlížení, zpracování, tvorby a~editace geodat, nad nimiž může provádět například SQL~dotazy. Data v~prostředí QGIS mohou být jak rastrová, tak také vektorová. Co se funkcionality týče, běžnému uživateli dostačuje, ale možností komerčních gigantů typu ArcGIS nedosahuje. 

Síla QGIS vedle volnosti užití tkví především ve značném množství veřejně přístupných zásuvných modulů~–~tzv.~pluginů. Ty byly zprvu psány především v~jazyku C++ (stejně jako základní tělo programu), nyní se čím dál častěji přechází k~jazyku Python, přičemž především pro grafické uživatelské rozhraní se využívá Qt~knihoven. 



\section{Qt Project}
\label{qt}

  \begin{figure}[H]
    \centering
      \includegraphics[width=120pt]{./pictures/qt.png}
      \caption[Qt logo]{Qt logo 
      (zdroj: \url{http://d3hp9ud7yvwzy0.cloudfront.net/wp-content/uploads/2015/02/Qt-logo-medium.png})}
      \label{fig:qt}
  \end{figure}

Qt představuje uživatelsky přívětivé multiplatformní vývojové prostředí s~důrazem na grafické
uživatelské rozhraní (\zk{GUI}). 

Vývoj Qt započali na úsvitu devadesátých let dva programátoři - Haavard Nord a~Eirik
Chambe-Eng. Původně dali své společnosti jméno Quasar technologies, ale brzy společnost
přejmenovali na Trolltech. Pod touto hlavičkou působili až do roku~2008, kdy firmu odkoupila
Nokia. V letech 2011 a~2012 pak probíhal další přesun, tentokráte pod podnik Digia. Digia
nadále spolupracuje s~Qt~Company. 

Qt není programovacím jazykem. Qt je framework psaný v~C++, umožňující ale práci i~v~dalších jazycích
včetně jazyka python (PyQt). Qt klade zesílený důraz na \zk{GUI} a~práci s~ním. Mezi jednotlivými
objekty se dá komunikovat a~vykonávat nad nimi operace pomocí tzv.~signálů a~slotů. \zk{GUI} se ukládá
do souboru s~příponou \textit{.ui}. Pro náš zásuvný modul je důležitá právě tato odnož Qt poskytovaná
samostatně jako Qt~Designer. 




